\begin{qfmechanism}[DeFi lending liquidations: fixed-spread and auction designs, close factors and health factors, and how much becomes liquidatable when prices fall]\label{mech:defi-lending-liquidation}
\qffield{Statement} Qin, Zhou, Gamito, Jovanovic and Gervais study Aave, Compound, MakerDAO and dYdX, over 85 percent of Ethereum lending, from their launch to April 2021. Two designs dominate: auctions (MakerDAO) and fixed-spread liquidation (Aave, Compound, dYdX), where a liquidator repays debt and claims collateral worth the repaid debt times $(1+LS)$, with LS the liquidation spread, up to a \emph{close factor}, the maximum share of the debt repayable in one liquidation (sec.~2.3, eq.~(1), PDF p.~3).

Over the period, liquidation proceeds were \$807.46 million in 28{,}138 events by 2{,}011 liquidator addresses; 641 auction liquidations were unprofitable for liquidators; Aave V2 had accumulated \$87.4 thousand of bad debt by the end of April 2021; and a 43 percent fall in ETH, like 13 March 2020, would make \$1.07 billion of collateral liquidatable on MakerDAO. Fixed-spread designs ``favor liquidators over borrowers'' because they allow more collateral than necessary to be sold (sec.~1, PDF p.~2) \citep[arXiv v2, secs.~1--2, PDF pp.~1--3]{roadmap-defi-liquidations}. Perez, Werner, Xu and Livshits find on Compound to September 2020 that ``variations of only 3\% in an asset's dollar price can result in over 10m USD becoming liquidable'', with over 70 percent of liquidatable positions liquidated immediately (abstract) \citep[arXiv v6, abstract]{oc-perez-knife-edge-2021}.

In Aave V3, the V2 rule that only half a position could be liquidated at once was relaxed: a position ``approaching insolvency (HF $<$ 0.95)'' can be liquidated in full; on layer-2 chains a grace period applies after a sequencer outage for $0.95<HF<1$ (PDF pp.~7--8 and~14) \citep[Aave V3 Technical Paper, 27 January 2022, PDF pp.~7--8 and~14]{al-aave}.

\qffield{Notation} HF the health factor, below one liquidatable; LS liquidation spread (bonus); CF close factor; LT liquidation threshold.

\qffield{Conditions} The figures are for 2019--2021 protocol versions and parameters, which governance changes; the HF definition in terms of collateral, liquidation thresholds and debt is the protocols' (not transcribed here).

\qffield{Practice} To assess a liquidation, record the protocol version and parameters at the pinned block, compute the maximum repay amount from the close factor and the HF band, and net the bonus of gas and DEX slippage before calling a liquidation profitable.

\qffield{Pitfalls} The Aave paper is the developers' description; parameters such as the e-mode example (97 percent LTV, 98 percent threshold, 2 percent bonus for stablecoins, PDF p.~6) are examples, not current values.

\qffield{See also} \hyperref[def:amm-pool-mechanics]{AMM pools}; \hyperref[def:crypto-trading-risk-layers]{crypto trading risk layers}; \hyperref[mod:liquidity-spirals]{liquidity spirals}.

\qffield{Further reading} \nolinkurl{oc-heimbach-huang-2024}; \nolinkurl{oc-lehar-parlour-2022}; \nolinkurl{oc-eskandari-oracles-2021}.

\end{qfmechanism}
